% Original: PDF strana 23, donji slajd.
\slajdid{02-s046}
\begin{frame}[label=02-s046]{Ista funkcija: drugačiji broj čipova}
\setlength{\parskip}{0pt}
\renewcommand{\arraystretch}{1.0}
Međutim, pri poređenju po broju čipova rezultat može biti drugačiji ili isti.
\par\vspace{1.5mm}
U oba oblika potrebna su četiri invertora za \(\bar D,\bar C,\bar B,\bar A\).
\par\vspace{1.5mm}
% element: 02-s046:e01
\begin{columns}[T,onlytextwidth]
\begin{column}{.48\textwidth}
\Oznaka{Zbir proizvoda — ukupno \textbf{4 čipa}}\Oznaka{\(F=\bar D\bar B+D\bar C\bar A\)}
\par\vspace{1.5mm}
\begin{tabular}{@{}ll@{}}
4 invertora & 1 čip \\
1 dvoulazno I & 1 čip \\
1 troulazno I & 1 čip \\
1 dvoulazno ILI & 1 čip \\
\end{tabular}

\end{column}
\begin{column}{.48\textwidth}
% element: 02-s046:e02
\Oznaka{Proizvod zbirova — ukupno \textbf{3 čipa}}\Oznaka{\(F=(D+\bar B)(\bar D+\bar C)(\bar D+\bar A)\)}
\par\vspace{1.5mm}
\begin{tabular}{@{}ll@{}}
4 invertora & 1 čip \\
3 dvoulazna ILI & 1 čip \\
1 troulazno I & 1 čip \\
\end{tabular}

\end{column}
\end{columns}
\par\vspace{1.5mm}
U 14-pinskom kućištu dva pina služe za masu i napajanje, pa ostaje 12 signalnih pinova:
\par
Invertor: 1 ulaz + 1 izlaz \(\Rightarrow (14-2)/2=6\) invertora u jednom čipu.
\par
Dvoulazno ILI: 2 ulaza + 1 izlaz \(\Rightarrow (14-2)/3=4\) kola u jednom čipu.
\par\vspace{1.5mm}
% element: 02-s046:e03
\alert{\textbf{Prema broju čipova povoljniji je proizvod zbirova}}, iako ima više ulaza kola.
\end{frame}
